\SourceLeaf{275}
\[
\begin{aligned}
Q+q&=46^0\ 40'\ 37''69794\\
Q+q&=58\quad39\quad44{,}09285
\end{aligned}
\]
die endliche Formel hingegen
\[
\begin{aligned}
Q+q&=46^0\ 40'\ 37''69794\\
Q+q&=58\quad39\quad44{,}09283
\end{aligned}
\]
also so genau übereinstimmend, wie zehnziffrige Logarithmen nur verstatten.
\Article{7}
\indent Auf ähnliche Weise lässt sich der Logarithm von \(m\) in eine Reihe entwickeln, deren erste Glieder folgende sind:
\[
\begin{aligned}
\log m={}&-\frac{\sin2\varphi^2}{6\cos\theta^4}\mathord{.}csp^3-\frac{\sin2\varphi^2}{24\cos\theta^6}(cc+11eess)p^4\\
&+\frac{\sin2\varphi^2}{120\cos\theta^8}\mathord{.}\frac{s}{c}(2cc-3ss-ee(40c^4-20ccss-6s^4)\\
&\qquad-e^4ss(104c^4+22ccss+3s^4))p^5
\end{aligned}
\]
Auch das folgende Glied habe ich (auf einem andern Wege) entwickelt, jedoch nur nach dem Hauptbestandtheile des Coëfficienten, welcher von der Ordnung \(ee\) ist, und dafür gefunden:
\[
+\frac{\sin2\varphi^2}{720\cos\theta^{10}}\mathord{.}\frac{1}{cc}(2c^4-18ccss-15s^4)p^6
\]
\indent Der durch diese Reihe ausgedrückte Logarithm ist der hyperbolische, und \(p\) wird, wie oben, in Theilen des Halbmessers ausgedrückt verstanden: verlangt man den briggischen Logarithmen, indem man \(p\) Grade bedeuten lässt, so muss noch der Modulus als Factor hinzukommen und \(\tfrac{\pi p}{180}\) für \(p\) geschrieben werden. In dieser Gestalt wird für unser Beispiel
\[
\begin{aligned}
\log m={}&-0{,}0049612433(\frac{p}{100})^3\\
&-0{,}0017329876(\frac{p}{100})^4\\
&-0{,}002393772(\frac{p}{100})^5\\
&-0{,}0124746(\frac{p}{100})^6
\end{aligned}
\]
